\documentclass[a4paper,12pt]{article}
\usepackage[T2A]{fontenc}
\usepackage[utf8]{inputenc}
\usepackage[russian]{babel}
\usepackage{amsmath,amssymb,mathtools}
\usepackage{geometry}
\geometry{left=23mm,right=22mm,top=23mm,bottom=23mm,headheight=26pt,headsep=6mm}
\usepackage{microtype}
\usepackage{tikz}
\usetikzlibrary{arrows.meta,positioning,shapes.geometric,calc}
\usepackage{xcolor}
\usepackage{titlesec}
\usepackage{enumitem}
\usepackage{booktabs,tabularx}
\usepackage{fancyhdr}
\usepackage{parskip}
\usepackage{icomma}
\usepackage{hyperref}
\definecolor{accent}{HTML}{008C91}
\definecolor{darkaccent}{HTML}{00545A}
\definecolor{textgray}{HTML}{404449}
\hypersetup{hidelinks,pdftitle={Чек-лист. Геометрический смысл производной},pdfauthor={Лёвин Артём Александрович}}
\setlength{\parindent}{0pt}
\setlength{\parskip}{7pt}
\linespread{1.13}
\setlist[enumerate]{leftmargin=*,itemsep=11pt,topsep=8pt}
\titleformat{\section}{\large\color{darkaccent}}{\thesection.}{0.65em}{}[\vspace{4pt}\color{accent}\titlerule]
\titlespacing*{\section}{0pt}{21pt}{12pt}
\titleformat{\subsection}{\normalsize\color{darkaccent}}{}{0pt}{}
\titlespacing*{\subsection}{0pt}{11pt}{3pt}
\pagestyle{fancy}
\fancyhf{}
\fancyhead[L]{\small\color{textgray}Геометрический смысл производной}
\fancyhead[R]{\small\color{textgray}Матвей Горбачёв}
\fancyfoot[C]{\small\thepage}
\renewcommand{\headrulewidth}{0.45pt}
\renewcommand{\headrule}{\hbox to\headwidth{\color{accent}\leaders\hrule height \headrulewidth\hfill}}
\newcommand{\key}[1]{\textcolor{darkaccent}{#1}}
\newcommand{\sep}{\,\text{; }\,}
\begin{document}
\begin{titlepage}
\thispagestyle{empty}
\centering
\vspace*{23mm}
{\large\color{darkaccent}МАТЕМАТИКА\par}
\vspace{20mm}
{\LARGE Чек-лист\par}
\vspace{16mm}
{\LARGE\color{accent}Геометрический смысл\\[3pt] производной и касательная\par}
\vspace{16mm}
{\normalsize Подготовка к ЕГЭ по профильной математике\par}
\vfill
{\large Матвей Горбачёв\par}
\vspace{10mm}
{\normalsize 1 ноября 2026 г.\par}
\vspace{14mm}
{\small Лёвин Артём Александрович\\[-1pt]
эксклюзивно для Матвея Горбачёва\par}
\begin{tikzpicture}[remember picture,overlay]
\draw[accent,line width=0.95pt]
([xshift=-4.5cm,yshift=24mm]current page.south)
.. controls +(+1.1cm,+.55cm) and +(-1.7cm,-.28cm) .. ([xshift=-.1cm,yshift=24mm]current page.south)
.. controls +(+1.7cm,+.39cm) and +(-1.5cm,-.37cm) .. ([xshift=4.5cm,yshift=24mm]current page.south);
\fill[darkaccent] ([xshift=-.1cm,yshift=24mm]current page.south) circle (1.25pt);
\end{tikzpicture}
\end{titlepage}
\setcounter{page}{1}

\section{Главная идея}
В каждой точке, где у функции существует производная, можно провести касательную к её графику. \key{Производная показывает наклон касательной} и позволяет численно описать, насколько быстро функция изменяется вблизи этой точки.

Для функции $y=f(x)$ в точке $x_0$ справедливо:
\[
\boxed{\displaystyle f'(x_0)=k=\tg\alpha},
\]
где $k$ --- угловой коэффициент касательной, а $\alpha$ --- угол её наклона к положительному направлению горизонтальной оси (если касательная не вертикальна).

\section{Как найти производную по чертежу}
Если на касательной известны две разные точки $A(x_1;y_1)$ и $B(x_2;y_2)$, то
\[
\boxed{\displaystyle f'(x_0)=k=\frac{y_2-y_1}{x_2-x_1}},\qquad x_1\ne x_2.
\]
Это отношение изменения по вертикали к изменению по горизонтали. Для вычисления удобно брать \key{узловые точки} --- точки с целочисленными координатами.

\begin{center}
\begin{tikzpicture}[scale=.83,>=Stealth]
\draw[step=1,gray!20,very thin] (-.3,-.35) grid (6.1,4.7);
\draw[->,gray!75] (-.35,0)--(6.4,0) node[right] {$x$};
\draw[->,gray!75] (0,-.45)--(0,4.95) node[above] {$y$};
\draw[accent,line width=1.1pt] (.3,.65)--(6.05,3.525);
\draw[densely dashed, darkaccent] (1,1)--(5,1)--(5,3);
\fill[darkaccent] (1,1) circle (2.4pt) node[below left] {$A(1;1)$};
\fill[darkaccent] (5,3) circle (2.4pt) node[above] {$B(5;3)$};
\node[below] at (3,1) {$\Delta x=4$};
\node[right] at (5,2) {$\Delta y=2$};
\draw[darkaccent] (4.64,1)--(4.64,1.36)--(5,1.36);
\end{tikzpicture}
\end{center}
\[
 f'(x_0)=\frac{3-1}{5-1}=\frac24=\frac12.
\]

\section{Знак производной}
\begin{center}
\begin{tabularx}{\linewidth}{@{}>{\raggedright\arraybackslash}p{.24\linewidth}X@{}}
\toprule
\textbf{Условие} & \textbf{Что показывает наклон касательной}\\
\midrule
$f'(x_0)>0$ & Касательная поднимается слева направо.\\
$f'(x_0)<0$ & Касательная опускается слева направо.\\
$f'(x_0)=0$ & Касательная горизонтальна.\\
\bottomrule
\end{tabularx}
\end{center}
\key{Проверка знака обязательна.} Например, если касательная идёт вниз, а длины катетов равны $3$ и $6$, то $f'(x_0)=-\frac36=-\frac12$.

\section{Касательная и уравнение прямой}
Для любой невертикальной прямой $y=kx+b$ число $k$ определяет наклон. Если эта прямая касается графика $f$ в точке с абсциссой $x_0$, то
\[
\boxed{f'(x_0)=k.}
\]
Если две невертикальные прямые параллельны, их угловые коэффициенты равны.

\subsection{Пример 1. Касательная задана уравнением}
Касательная к графику $f$ в точке $x_0$ имеет уравнение $y=-3x+7$. Тогда $f'(x_0)=-3$. Свободный член $7$ на величину производной не влияет.

\subsection{Пример 2. Касательная параллельна прямой}
Касательная параллельна прямой $y=5x-2$. Следовательно, $f'(x_0)=5$.

Если касательная параллельна $y=-20$, она горизонтальна, и $f'(x_0)=0$.

\section{График функции и график производной}
Перед решением графической задачи установите, \key{какой именно график дан}: $y=f(x)$ или $y=f'(x)$. Смысл чтения рисунка различается.

\begin{center}
\begin{tabularx}{\linewidth}{@{}>{\raggedright\arraybackslash}p{.25\linewidth}X@{}}
\toprule
\textbf{Дан график} & \textbf{Что искать}\\
\midrule
$f(x)$ & На участках возрастания $f'>0$, на участках убывания $f'<0$, если производная существует.\\
$f'(x)$ & Выше оси $Ox$: $f$ возрастает; ниже оси $Ox$: $f$ убывает. Нули $f'$ --- возможные точки экстремума.\\
\bottomrule
\end{tabularx}
\end{center}

\key{Тонкость.} Из $f'(x_0)=0$ ещё не следует, что $x_0$ --- точка экстремума. Например, для $f(x)=x^3$ в нуле производная равна нулю, но функция продолжает возрастать. Если знак производной меняется с $+$ на $-$, возникает максимум; с $-$ на $+$ --- минимум.

\subsection{Пример 3. Анализ знаков}
Пусть $f'(x)=x^2-4=(x-2)(x+2)$. Тогда $f'(x)>0$ при $x<-2$ и $x>2$, а $f'(x)<0$ при $-2<x<2$. Поэтому $f$ возрастает на $(-\infty;-2)$ и $(2;+\infty)$, убывает на $(-2;2)$, имеет максимум при $x=-2$ и минимум при $x=2$.

\section{Точка касания: два обязательных условия}
Пусть графика функции $y=f(x)$ касается прямая $y=g(x)=kx+b$ в точке с абсциссой $x_0$. В этой точке совпадают и значения функций, и их наклоны:
\[
\boxed{\begin{cases}
f'(x_0)=k,\\
f(x_0)=kx_0+b.
\end{cases}}
\]
Первое уравнение определяет \key{кандидатов} на точку касания. Если кандидатов несколько, второй строкой системы проверяют, какие из них действительно лежат на обеих кривых.

\subsection{Пример 4. Одна возможная абсцисса}
Прямая $y=6x+b$ касается параболы $f(x)=x^2+7x+1$. Найдём абсциссу касания и $b$.
\[
f'(x)=2x+7,\qquad 2x_0+7=6\quad\Longrightarrow\quad x_0=-\frac12.
\]
Затем сравним значения функций в этой точке:
\[
f\left(-\frac12\right)=\frac14-\frac72+1=-\frac94,
\qquad 6\left(-\frac12\right)+b=-3+b.
\]
Откуда $-3+b=-\frac94$, то есть $\boxed{b=\frac34}$.

\subsection{Пример 5. Несколько кандидатов}
Проверим касание графика $f(x)=x^3-2x^2-6x-4$ прямой $y=-2x-12$.
\[
f'(x)=3x^2-4x-6,\quad f'(x)=-2
\quad\Longrightarrow\quad 3x^2-4x-4=0.
\]
Получаем $x_1=2$ и $x_2=-\frac23$. Для $x=2$:
\[
f(2)=-16,\qquad -2\cdot2-12=-16.
\]
Значения совпали. При $x=-\frac23$ значения различаются:
\[
f\left(-\frac23\right)=-\frac{32}{27},\qquad
-2\left(-\frac23\right)-12=-\frac{32}{3}.
\]
Поэтому прямая касается графика только при $\boxed{x_0=2}$.

\section{Как найти параметр из условия касания}
Иногда неизвестны и абсцисса $x_0$, и параметр функции. Тогда используйте обе строки системы и решайте её как обычную алгебраическую систему.

\subsection{Пример 6. Неизвестный коэффициент}
Прямая $y=x-3$ касается параболы $f(x)=ax^2+3x-2$. Найдём $a$.
\[
\begin{cases}
2ax_0+3=1,\\
ax_0^2+3x_0-2=x_0-3.
\end{cases}
\quad\Longrightarrow\quad
\begin{cases}
ax_0=-1,\\
ax_0^2+2x_0+1=0.
\end{cases}
\]
Подставляем $ax_0^2=(ax_0)x_0=-x_0$ во вторую строку:
\[
-x_0+2x_0+1=0\quad\Longrightarrow\quad x_0=-1.
\]
Из $ax_0=-1$ получаем $\boxed{a=1}$.

\clearpage
\thispagestyle{fancy}
\begin{center}
{\large\color{darkaccent}Алгоритм: проверка касания прямой и графика}\par
\vspace{16mm}
\begin{tikzpicture}[
 node distance=13mm,>=Stealth,
 action/.style={draw=darkaccent,rounded corners=2pt,align=center,text width=10.3cm,minimum height=13mm,inner sep=7pt,font=\normalsize},
 terminal/.style={draw=accent,ellipse,align=center,minimum width=5.1cm,minimum height=11mm,font=\normalsize},
 decision/.style={draw=darkaccent,diamond,aspect=3.2,align=center,text width=6.7cm,inner sep=5pt,font=\normalsize},
 arr/.style={-Stealth,thick,draw=darkaccent}
]
\node[terminal] (start) {Дана $f(x)$ и прямая $g(x)=kx+b$};
\node[action,below=of start] (der) {Вычислить $f'(x)$ и записать уравнение $f'(x)=k$};
\node[action,below=of der] (solve) {Найти все действительные решения $x$ первого уравнения};
\node[action,below=of solve] (verify) {Для каждого кандидата проверить второе условие: $f(x)=kx+b$};
\node[decision,below=15mm of verify] (has) {Остались кандидаты, удовлетворяющие обоим условиям?};
\node[terminal,below=20mm of has] (yes) {Записать найденные точки касания};
\node[action,right=13mm of has,text width=3.4cm,minimum height=14mm,font=\small] (no) {Для заданной прямой касания нет};
\draw[arr] (start)--(der);
\draw[arr] (der)--(solve);
\draw[arr] (solve)--(verify);
\draw[arr] (verify)--(has);
\draw[arr] (has)--node[left,font=\small]{да}(yes);
\draw[arr] (has.east)--node[above,font=\small]{нет}(no.west);
\end{tikzpicture}
\end{center}
\clearpage

\section*{Тренировочный блок \normalsize(Путь А)}
\vspace{-5mm}
Решите 10 заданий по возрастанию сложности. В задачах на касательную сначала запишите равенство производных, затем при необходимости проверьте совпадение значений функций.

\subsection{Базовый уровень}
\begin{enumerate}[label=\textbf{\arabic*.},itemsep=8pt]
\item Касательная к графику $f(x)$ проходит через $A(1;2)$ и $B(5;10)$. Найдите $f'(x_0)$ в точке касания.
\item Касательная проходит через $A(-2;6)$ и $B(4;-3)$. Найдите производную в точке касания.
\item Касательная к графику $f(x)$ проходит через начало координат и точку $M(3;-6)$. Найдите $f'(x_0)$.
\end{enumerate}

\subsection{Средний уровень}
\begin{enumerate}[label=\textbf{\arabic*.},start=4,itemsep=8pt]
\item Касательная к графику $f(x)$ в точке $x_0$ параллельна прямой $y=5x-7$. Найдите $f'(x_0)$.
\item В точке $x_1$ касательная проходит через $A(0;2)$ и $B(4;4)$, а в точке $x_2$ --- через $C(-2;5)$ и $D(2;1)$. Найдите $f'(x_1)-f'(x_2)$.
\item Известно, что $f'(x)=x^2-9$. Укажите интервалы возрастания и убывания функции и абсциссы её точек максимума и минимума.
\item К графику $f(x)=x^2+5x+2$ проведена касательная, параллельная прямой $y=3x-4$. Найдите абсциссу точки касания.
\end{enumerate}

\subsection{Повышенный уровень}
\begin{enumerate}[label=\textbf{\arabic*.},start=8,itemsep=8pt]
\item Прямая $y=2x+b$ касается графика $f(x)=x^2+4x+1$. Найдите $b$.
\item Прямая $y=2$ касается графика $f(x)=x^3-3x$. Найдите абсциссу точки касания.
\item Прямая $y=x-4$ касается графика $f(x)=ax^2+5x-1$. Найдите коэффициент $a$.
\end{enumerate}

\clearpage
\section*{Ответы}
\renewcommand{\arraystretch}{1.46}
\begin{center}
\begin{tabularx}{\linewidth}{@{}>{\centering\arraybackslash}p{1.1cm}X@{}}
\toprule
\textbf{№} & \textbf{Ответ}\\
\midrule
1 & $2$\\
2 & $-\dfrac32$\\
3 & $-2$\\
4 & $5$\\
5 & $\dfrac32$\\
6 & Возрастает на $(-\infty;-3)$ и $(3;+\infty)$, убывает на $(-3;3)$; максимум при $x=-3$, минимум при $x=3$.\\
7 & $-1$\\
8 & $0$\\
9 & $-1$\\
10 & $\dfrac43$\\
\bottomrule
\end{tabularx}
\end{center}

\vspace{12mm}
\subsection{Итоговая памятка}
\begin{center}
\begin{tabularx}{\linewidth}{@{}>{\raggedright\arraybackslash}p{.36\linewidth}X@{}}
\toprule
\textbf{Что известно} & \textbf{Что делать}\\
\midrule
Две точки касательной & Найти $k=\dfrac{\Delta y}{\Delta x}$ и проверить знак.\\
Уравнение $y=kx+b$ & Взять $f'(x_0)=k$.\\
Касательная параллельна прямой & Приравнять угловые коэффициенты.\\
Дан график $f'$ & Использовать знак $f'$ для определения поведения $f$.\\
Нужно найти точку касания & Решить $f'(x_0)=k$ и проверить $f(x_0)=kx_0+b$.\\
Неизвестный параметр & Решить систему из двух условий касания.\\
\bottomrule
\end{tabularx}
\end{center}
\end{document}
